\chapter{Que date la date ? Événements, séquences et gravure}\label{ch:evenements}

Une année consulaire n’est intelligible qu’avec l’événement auquel elle se rapporte. La mention d’une naissance ne date pas une mort ; l’achat d’un emplacement ne date pas la déposition ; une commémoration collective peut juxtaposer des événements éloignés. L’année de classement du catalogue est donc un instrument de repérage, non une variable historique immédiatement univoque.

\section{Identifier l’événement avant de classer la date}

Le relevé associe à chaque indication calendaire un événement, une personne lorsque celle-ci est identifiable, un statut textuel et une justification. Les verbes de décès (\emph{decessit, recessit, defunctus, quievit}) et les termes de déposition (\emph{depositus, depositio}) fournissent des indices distincts. L’association d’une date au mot le plus proche aide au repérage ; elle ne remplace pas l’analyse de la syntaxe et de la structure du document. En l’absence de relation explicite, la catégorie reste « événement indéterminé » : c’est le cas de \st{ev_indet} dates sur \st{nDatesConv} dans l’ensemble C.

\section{Décès et déposition}

Les \nnums{119} et 211 distinguent un décès et une mise au tombeau. Leur intérêt est d’attester que deux opérations peuvent recevoir deux dates dans le même document. Ils ne suffisent pas à fournir un délai habituel d’inhumation, ni à expliquer tous les écarts entre quantième et jour de la semaine. Un écart peut relever de la lecture, du calcul, d’une notation inclusive ou de la relation entre événements ; ces possibilités doivent rester distinctes tant qu’un témoin supplémentaire ne permet pas de décider.

\section{Documents à plusieurs indications chronologiques}

Les \nnums{15}, 55 et 126 obligent à séparer une naissance d’un événement funéraire. D’autres documents juxtaposent plusieurs défunts. La grande dalle du \nnum{134} (105 × 190 cm selon l’EDB) porte plusieurs textes : la \emph{dormitio} du prêtre Dulcitius, celle du diacre Gregorius et celle de leur sœur, puis la déposition de Faustina, datée par le consulat de Valentinien \emph{nobilissimus puer} et de Victor (369), et la \emph{dormitio} d’Agapè. L’EDB date l’ensemble « 350–399 ; 369 » : l’année de classement n’est portée que par l’une des parties, et les autres dates ne sont pas nécessairement de la même année. La plaque ICUR IV 11117 porte de même deux dépositions successives, d’un défunt en 389 (\nnum{313}) et de son frère (\nnum{314}), dont l’année est restituée. Les \nnums{117}, 236, 317 et 373 demandent une analyse semblable de la succession des commémorations. L’attribution de chaque date à un défunt ou à une intervention doit être contrôlée dans la notice, sans la déduire du seul nombre de consulats.

Pour les analyses du volume, la règle de classement reste l’année retenue dans le catalogue. Cette variable est expressément nommée « année de classement ». Une analyse de la mort doit employer les événements de décès ; une analyse des dépositions, les événements de déposition. Une étude de la production épigraphique doit distinguer les phases de gravure et ne peut prendre automatiquement la première ou la dernière date funéraire pour l’année de fabrication du support.

\section{Préparation du vivant et additions}

Une formule \emph{se vivo}, \emph{se viva} ou \emph{se vivi} signale une initiative prise du vivant d’un commanditaire. Elle ne garantit pas que toute la surface conservée a été gravée au même moment. Le \nnum{300} associe ainsi une préparation du vivant et une mention de décès ; l’ordre exact des opérations ne peut être établi sur la seule continuité du texte transmis.

La date du dernier événement mentionné fournit au plus tôt une limite pour l’état complet du texte, sous réserve de l’interprétation des événements et de l’absence d’addition postérieure. Elle ne suffit pas à dater toutes les lignes ni le premier emploi du support. La matérialité de l’inscription est ici une condition de l’analyse chronologique, non un complément décoratif.

\section{Quand la date n’appartient qu’à un témoin}

La confrontation des bases ajoute un cas de figure. Pour huit textes, la formule de date n’est transmise que par l’EDB (chapitre~\ref{ch:fondement}). Inversement, au \nnum{378}, la formule \emph{[Olybrio et Pro]bino} (395) et la généalogie cléricale qui l’accompagne ne figurent dans aucune des fiches EDB que l’on peut rattacher au monument (ICVR III 8345, 8274c et 8387v) : l’EDB date l’ensemble « 390–424 ». L’année de classement repose alors sur une partie du texte qu’un seul témoin transmet. Elle est conservée, mais elle ne peut servir de jalon tant que l’ICVR III 8345 n’a pas été contrôlée.

\section{Faces, champs et remploi}

Une rupture de champ dans un export n’est pas une preuve de remploi. Elle peut représenter deux faces, deux zones d’une même plaque, une addition ou une convention de saisie. Les fiches de l’EDB précisent parfois la disposition : la colonne de Saint-Paul (\nnum{321}) porte trois textes en trois endroits du fût ; la dalle du \nnum{134} a été remployée en morceaux dans un mur. Ailleurs, la qualification d’un support comme opisthographe ou remployé exige une édition descriptive, une photographie ou un examen matériel.
